
\section{Meromorphic and Analytic Continuation}\label{sec:meromorphic-analytic}

The analytic continuation of the Riemann zeta function, and more generally of Dirichlet $L$-functions, requires the notion of a meromorphic function. We record the relevant definitions and a motivating example here.

\begin{Definition}{Zero, Pole, Order, Residue}{}
  \index{zero!of a function}\index{pole}\index{residue}%
  Assume $f$ has a convergent Laurent series expansion about $z_0$:
  $$f(z) = a_n(z-z_0)^n + a_{n+1}(z-z_0)^{n+1} + \cdots = \sum_{m=n}^{\infty} a_m(z-z_0)^m,$$
  with $a_n \neq 0$. If $n > 0$ we say $f$ has a \emph{zero of order $n$} at $z_0$; if $n < 0$ we say $f$ has a \emph{pole of order $-n$} at $z_0$. If $n = -1$ we say $f$ has a \emph{simple pole with residue $a_{-1}$}. We denote the order of $f$ at $z_0$ by $\ord_f(z_0) = n$.
\end{Definition}

\begin{Definition}{Meromorphic Function}{}
  \index{meromorphic function}\index{analytic function}%
  We say $f$ is \emph{meromorphic at $z_0$} if there exists a disk about $z_0$ of radius $r$ and an integer $n_0$ such that for all $z$ with $|z - z_0| < r$,
  $$f(z) = \sum_{n \geq n_0} a_n(z-z_0)^n.$$
  If $f$ is meromorphic at each point in a region, we say $f$ is \emph{meromorphic} on that region. If $n_0 \geq 0$, we say $f$ is \textbf{analytic}.
\end{Definition}

A canonical example illustrates the notion of continuation. Consider the geometric series $G(r) = \sum_{n=1}^{\infty} r^n$, which converges for $|r| < 1$ and equals $\frac{1}{1-r}$ there. Define $H(r) = \frac{1}{1-r}$; then $H$ is well defined for all $r \neq 1$ and agrees with $G$ wherever $G$ converges. Since $H$ is defined on a strictly larger domain and has a simple pole at $r = 1$ (with residue $-1$), we call $H$ a \textbf{meromorphic continuation}\index{meromorphic continuation} of $G$. Had $H$ no poles, we would call it an \textbf{analytic continuation}\index{analytic continuation}. A function that is analytic at every point of $\bbC$ --- equivalently, one whose Taylor expansion converges everywhere, so that it has no poles anywhere --- is said to be \textbf{entire}\index{entire function}.

\section{Fermat's Little Theorem}
Fermat's Little Theorem is one of the oldest results in elementary number theory, recorded in a 1640 letter from Pierre de Fermat to Fr\'{e}nicle de Bessy. It underlies essentially every computation we perform in the multiplicative group $\bbF_p^*$ -- from the very existence of the canonical additive character of $\bbF_q$ (where $q=p^r$), through the character congruence $\psi(l)\equiv l^f\pmod P$ used in the arithmetic proof of the Davenport--Hasse relation, down to the mod-$q$ reduction of Gauss sums appearing in the proof of the \QRthm. We record it here for completeness.
\begin{Theorem}{Fermat's Little Theorem}{thm:fermat-little}%
  \index{Fermat's little theorem}
  Let $p$ be a prime. Then for every integer $a$ we have $$a^p\equiv a\pmod p,$$ and in particular, if $p\nmid a$ then $a^{p-1}\equiv 1\pmod p$.
\end{Theorem}
\begin{proof}
  The non-zero residues modulo $p$ form a group $\bbF_p^*=\{1,2,\dots,p-1\}$ under multiplication of order $p-1$. By Lagrange's theorem every element of a finite group satisfies $g^{|G|}=e$, so $a^{p-1}\equiv 1\pmod p$ for every $a\in\bbF_p^*$, i.e.\ every $a$ with $p\nmid a$. Multiplying both sides by $a$ gives $a^p\equiv a\pmod p$ for such $a$; and the congruence $a^p\equiv a\pmod p$ is trivially true when $p\mid a$, so it holds for every integer $a$.
\end{proof}
